\documentclass[12pt,a4paper]{report}
\usepackage[]{amsmath,amsthm,amssymb,amscd}
\usepackage[latin1]{inputenc}
%\usepackage{fullpage}

\topmargin 0pt
\advance \topmargin by -\headheight
\advance \topmargin by -\headsep
     
\textheight 246mm
     
\oddsidemargin 0pt
\evensidemargin \oddsidemargin
\marginparwidth 0.5in
     
\textwidth 159mm


%               Theorem Environments
\theoremstyle{definition}
\newtheorem{ex}{}
\newcommand{\pd}[1]{\frac{\partial}{\partial #1}} %\pd{x}
%               Subquestions numbered with letters
\renewcommand{\theenumi}{\textbf{\alph{enumi}}}

%               Maths definitions

\newcommand{\NN}{\ensuremath{\mathbb N}}
\newcommand{\ZZ}{\ensuremath{\mathbb Z}}
\newcommand{\CC}{\ensuremath{\mathbb C}}
\newcommand{\TT}{\ensuremath{\mathbb T}}
\newcommand{\RR}{\ensuremath{\mathbb R}}
\newcommand{\QQ}{\ensuremath{\mathbb Q}}
\newcommand{\g}{\ensuremath{\frak{g}}}
\newcommand{\h}{\ensuremath{\frak{h}}}
\newcommand{\ft}{\ensuremath{\frak{t}}}
\newcommand{\cB}{\mathcal{B}}
\newcommand{\cN}{\mathcal{N}}
\newcommand{\cL}{\mathcal{L}}
\newcommand{\cD}{\mathcal{D}}
%\newcommand{\pd}[1]{\frac{\partial}{\partial #1}} %\pd{x}
%\newcommand{\pd}[1]{\frac{\partial}{\partial #1}} %\pd{x}

\begin{document}
\noindent
{\sffamily\large Marco Zambon 
\hfill     27/03/2012\\Geometr\'ia III, curso 2011-12}
 

%\noindent
%\hrulefill{}\\
\begin{center}\bfseries\large\textsf{Examen Parcial    \vspace{-2mm}}
\end{center}

 
%\hrulefill{}\\

\noindent
\hrulefill{}\\

\noindent
{\sffamily\large Tu nombre:
\\
Tu DNI: 

}


 
\begin{center}{\sffamily (deja esta tabla vac\'ia)}
\end{center}
\vspace*{-5mm}
%\begin{table}[htdp]
%\caption{}
%\begin{center}
%\begin{tabular}{|c|c|c|c|c|c||c|}
%\hline
%Ej. 1  &Ej. 2&Ej. 3&Ej. 4&Ej. 5&Ej. 6& Total puntos (de 100)\\
%\hline
%& & & & & &  \\
%& & & & & &  \\
%\hline
%\end{tabular}
%\end{center}
%\label{default}
%\end{table}
 

\begin{table}[htdp]
%\caption{}
\begin{center}
\begin{tabular}{|c|p{2cm} |}
\hline
&\\[-0.2cm]
Ej. 1  &\\[0.2cm]
\hline
&\\[-0.2cm]
Ej. 2  &\\[0.2cm]
\hline
&\\[-0.2cm]
Ej. 3  &\\[0.2cm]
\hline
&\\[-0.2cm]
Ej. 4  &\\[0.2cm]
\hline\hline
&\\[-0.2cm]
Total (de 30) & \\[0.2cm]
\hline
\end{tabular}
\end{center}
\label{default}
\end{table}
\vspace*{-5mm}


%%
\noindent
%\hrulefill{}\\
\textsf{{\bf \sffamily  Instrucciones:} Utiliza una hoja distinta para 
cada problema, y en cada hoja escribe  tu nombre y el n\'umero del problema.\\
Justifica todas tus respuestas.
Puedes utilizar, cit\'andolos de manera apropriada, resultados que hemos obtenido en clase o en los ejercicios para casa.    \vspace{-2mm}}

 
\noindent
\hrulefill{}\\

 
\begin{ex}[8 puntos]    
Considera $S^2=\{v\in \RR^3: |v|=1\}$  con la relaci\'on de equivalencia $$v\sim_a w \Leftrightarrow (v=w \text{ o bien } v=-w).$$
Considera tambi\'en  $\RR^3-\{0\}$ con la relaci\'on de equivalencia $$v\sim_b w \Leftrightarrow \exists \lambda \neq 0: v=\lambda \cdot  w.$$
Encuentra \emph{explicitamente} un homeomorfismo $$S^2/\sim_a \to (\RR^3-\{0\})/\sim_b,$$ y demuestra que es un homeomorfismo.

\noindent{\bf Nota:} En clase vimos que existe un tal homeomorfismo
(ambos cocientes son homeomorfos al plano proyectivo). Aqu\'i os pido de escribir \emph{explicitamente} un tal homeomorfismo.
\end{ex}

\begin{ex} [7 puntos]
Considera $$S:=([0,1]\times[0,1])-\cup_{i=1}^5 D_i,$$ donde  $D_1,\dots,D_5$ son  discos abiertos en  $(0,1)\times(0,1)$ tal que  $\bar{D_i}\cap \bar{D_j}=\emptyset$ si $i\neq j$. Considera 
$S/\sim$, donde la relaci\'on de equivalencia $\sim$ es la siguiente:  cada punto es equivalente a s\'i mismo, y adem\'as $(0,y)\sim (1,1-y)$ para todo $y\in [0,1]$. 

> A qu\'e superficie en la lista de superficies compactas, conexas, con borde es homeomorfa  $S/\sim$? Demuestralo.

\end{ex}
 

\begin{ex}  [8 puntos]
Sea $S^2=\{v\in \RR^3: |v|=1\}$, con su estructura   diferencial standard. Sea $\textbf{S}:=(0,0,-1)\in S^2$.
Para todo $(x,y,z)\in \RR^3$, se denota por $|(x,y,z)|:=\sqrt{x^2+y^2+z^2}$ su norma.

> Es la funci\'on
$$f \colon S^2\to \RR,  \;\;\; v\mapsto |v-\textbf{S}|$$
diferenciable? Demuestralo.

\noindent{\bf Nota:} Puedes utilizar que una carta diferencial de $S^2$ es dada por la proyecci\'on estereogr\'afica $\phi_N$, y que su inversa es $$(\phi_N)^{-1} \colon \RR^2 \to S^2-\{\textbf{N}\},\;\;(u,v)\mapsto \frac{1}{u^2+v^2+1}(2u,2v,u^2+v^2-1),$$
donde $\textbf{N}:=(0,0,1)$.

 
\end{ex}





\begin{ex}[7 puntos]
Si $M$ es una variedad diferencial y $S,R$ son subvariedades de $M$ tal que $dim(S)=dim(R)$, > es verdad que la union $S\cup R$ es siempre  una subvariedad de $M$? Demuestralo.
\end{ex}

 

%\begin{ex}  
%Demuestra que 
%$$T \;\;\sharp \;\; C\cong T-(D_1 \cup D_2)$$
%donde $T=S^1\times S^1$ es el toro, $C=S^1\times [0,1]$ el cilindro, y $D_i\subset T$ (por $i=1,2$) son discos cerrados tal que $D_1\cap D_2=\emptyset$.
%\end{ex} 






 \end{document}
 %%%%%%%
 \begin{ex}

\end{ex}

 \begin{ex}
\begin{enumerate}
\item 
\end{enumerate}
\end{ex}
 
 
 
 